% Source: sources/returns/S6_04_MAPS_MANIFOLDS/sections/manifolds.tex:2156–2362
\paragraph{Theorem D.2.11 (Proved: all orders $3k$ for $IV^*+II+I_2$).}
The rational elliptic surface
\[
 y^2=x^3-3t^3(t-1)x+t^4(t-1)(2t-1),
 \qquad P=(t^2,t^2),
 \tag{D.C1}
\]
has fibres $IV^*,II,I_2$ at $0,1,\infty$ and Mordell--Weil group
$\mathbb Z$ generated by $P$. Put $M=\mathcal O(P-O)$ and
$Q=3P$. Let $C_0$ be the identity component of the $I_2$ fibre.
There is a linearization with divisor $C_0$, and its torus
family has Yang's index-one Type C marking. Keep the type-$II$
filling unchanged and use the index-one toroidal filling at
$I_2$. At $IV^*$ choose any logarithmic order and invariant vector
\[
 m=3k,\qquad v=n(\gamma-u-w)+r\delta,
 \qquad k>0,\quad n,r\in\mathbb Z,
 \quad\gcd(n,3k)=1.
\]
The resulting compact smooth complex threefold has
\[
 \pi_1(X)=\mathbb Z/|n|,\qquad
 H_i(X;\mathbb Z)=
 \begin{cases}
 \mathbb Z&i=0,6,\\
 \mathbb Z/|n|&1\leq i\leq4,\\
 0&i=5.
 \end{cases}
\]
Consequently $n=\pm1$ gives the standard $S^6$ for every positive
$k$. This is an affirmative unbounded configuration for
Viaclovsky's Question 4. The linearization is supported on one
$I_2$ component, so the local-narrowness restriction of the public
engine does not include this input.
\begin{proof}
\emph{The explicit elliptic surface and linearization.}
For the coefficients $f,g$ of (D.C1), direct calculation gives
\[
 4f^3+27g^2=27t^8(t-1)^2,\qquad
 j=6912t(1-t).
\]
The minimal valuations give $IV^*$ at zero and $II$ at one.
In the chart at infinity the discriminant has order two and
the cubic has a node, giving $I_2$. These fibres have Euler
sum twelve; the degree-four and degree-six Weierstrass
coefficients give the rational elliptic surface. Substitution
checks the section $P$, which is disjoint from $O$ and meets
the nonidentity components at $IV^*$ and $I_2$.
Its height is $2-4/3-1/2=1/6$. The Shioda--Tate rank is one,
and OS entry 49 has height lattice $\langle1/6\rangle$ and
no torsion, as checked in D.2.1; hence $P$ is a generator.

The section $Q=3P$ meets the identity component of $IV^*$
and the nonidentity component of $I_2$. Its height is $3/2$,
so the height formula gives $Q\cdot O=0$.
The mixed height $\langle P,Q\rangle=1/2$, with the sole
local contribution $1/2$ at $I_2$, gives $P\cdot Q=0$.
Let $J=M\otimes t_Q^*M^{-1}$. Its multidegree vanishes at
$IV^*$ and $II$. At $I_2$, translation by $Q$ exchanges the
two components, so its multidegrees are $(-2,2)$ on
$(C_0,C_1)$, exactly those of $\mathcal O(C_0)$.
Their difference is thus a pullback from the base. Restriction
to $O$ gives
\[
 J\cdot O=Q\cdot O-P\cdot Q+1=1=C_0\cdot O,
\]
so the base degree of that difference is zero. Therefore
$J\simeq\mathcal O(C_0)$, supplying the asserted effective
linearization. Its norm under the component exchange has divisor
$C_0+C_1=F_\infty$. This explains why a half full-fibre weight
appears in period coordinates while the geometric divisor is
integral.

\emph{Identification with the actual Type C periods.}
The elliptic $j$-function is the one in
\cite[Proposition 3.1]{manYang26}. Over $t=s^3$ the good
model has central curve $Y^2=X^3+1$, and $P$ specializes to
$(0,1)$, a nonzero fixed point of order three.
Over $t-1=s^6$, the same section specializes to the origin
of the good elliptic model. At $I_2$ it meets the nonidentity
component. Choose the sign of the elliptic marking so that its
Abel--Jacobi coordinate at the first point is $(\tau+1)/3$.
Its normalized lift therefore obeys the two Type C affine
laws and has $\mu-\tau/2$ regular at infinity.
The uniqueness proof of \cite[Proposition 3.4]{manYang26},
using $H^0(\mathcal O(-1))=H^1(\mathcal O(-1))=0$, identifies
it with $\mu_C$.
The extension line $M$ has this same Picard coordinate,
whereas the translation has coordinate $3\mu_C$.
The norm divisor just calculated fixes the cusp central
period. The remaining scalar period is the affine constant
of \cite[Proposition 3.6]{manYang26}. Consequently the open
two-torus family is the index-one Type C family, for a suitable
admissible constant, and its primitive toroidal cusp has one
normalized $dP_6$ component.

An exact integral check of the marking is also available.
In the order $(II,IV^*,I_2)$, the homology matrices are
\[
 \begin{split}
 M_{II}&=\begin{pmatrix}1&1&0&0\\-1&0&0&0\\0&0&1&0\\0&0&0&1\end{pmatrix},\\
 M_{IV}&=\begin{pmatrix}0&-1&0&0\\1&-1&0&-3\\0&1&1&1\\0&0&0&1\end{pmatrix},\\
 M_{I_2}&=\begin{pmatrix}1&2&0&3\\0&1&0&0\\0&-1&1&-1\\0&0&0&1\end{pmatrix}.
 \end{split}
\]
The first two follow from the locally narrow finite-point
section cocycles $(0,0)$ and $(0,-1)$ for $P$, multiplied by
three for $Q$. Their product determines the third matrix;
no local-narrowness hypothesis at $I_2$ is used.
Put $U=M_{IV}$, $V=U^{-1}M_{II}U$ and
\[
 S=\begin{pmatrix}
 0&0&-1&0\\0&-1&1&1\\0&0&0&1\\-1&0&3&0
 \end{pmatrix}.
\]
Then $\det S=-1$, $U^{-t}S=ST_1$ and $V^{-t}S=ST_2$,
where $T_1,T_2$ are precisely Yang's displayed Type C
cohomological matrices. This is the finite Hurwitz reorder
followed by an integral change of basis, not merely a rational
comparison of eigenvalues.

\emph{The unchanged type-$II$ filling with its integral marking.}
Here the product proof of \cite[Lemma 3.10]{manLR26} applies
with $3P$ in place of its translation section; we check all the
needed ingredients. A closed elliptic neighborhood $E$ of the
type-$II$ fibre retracts to its cuspidal cubic, homeomorphic to
$S^2$. The restriction of $M$ has degree zero and is therefore
differentiably trivial. In a smooth trivialization the lifted
contracting generator has form
\[
 (x,z)\longmapsto(t_{-Q}(x),\lambda(x)z),\qquad |\lambda|<1.
\]
As $E$ is simply connected, $\lambda$ has a smooth logarithm.
The relative exponential of the additive N\'eron group is an
isomorphism on the central smooth locus. After shrinking the
disc it provides a holomorphic logarithm of $-Q$ and hence a
fibre-preserving isotopy of $t_{-Q}$ to the identity.
The translations extend over the cuspidal cubic and fix its
cusp. Combining this isotopy with the logarithm of $\lambda$
trivializes the circle-bundle mapping torus of a fundamental
annulus. Thus this filling is homeomorphic over its base disc
to $E\times T^2$, with angular circle $\delta$ and a primitive
second circle of $\gamma^*$-value one.

Its regular specialization sends the elliptic area $uw$ to
the fundamental class of the cuspidal cubic and $\gamma\delta$
to the area of this second torus; it kills the mixed classes.
Equivalently the exact degree-two pullback lattice is
$\langle u^*w^*,\gamma^*\delta^*\rangle$.
There is no index-six factor: this is the unchanged filling,
not a free order-six affine quotient.
Both transverse normal tori bound products of a section disc
with a circle. In particular its $\delta$ meridian sweep is
zero. These assertions also follow by applying the marked
specialization proof of \cite[Proposition 3.11]{manLR26}
to the product just constructed.

\emph{All orders and all displayed twists.}
At $IV^*$ the order-$3k$ free quotient and its marked collar
are those of D.2.10 with $R=1$, $n_1=n$, $r_1=r$.
The two other compact fillings have just been specified,
so the same finite-core gluing proves compactness and
smoothness. The unchanged type-$II$ piece kills $u,w$ and
its meridian; the cusp kills $\delta$ and the remaining
meridian. Van Kampen leaves exactly $n\gamma=0$, proving
$\pi_1(X)=\mathbb Z/|n|$.

The first finite pullback matrix is unchanged. For the second
piece the primitive product lattice replaces its coefficient
$m_2$ by one, and its vanishing $\delta$ sweep replaces $n_2$
by zero. Thus the complete augmented integral map is
\[
 D=\begin{pmatrix}
 -1&1&0&1&0&0&0\\
 k&-k&3k&0&0&0&-n\\
 0&0&0&-1&0&0&0\\
 0&0&1&0&0&0&0
 \end{pmatrix}.
\]
The primitive $w$ cusp sweep still kills the degree-one
Mayer--Vietoris kernel. Unit elimination reduces $D$ to
one unit factor and
\[
 \begin{pmatrix}1&0&0\\-k&3k&-n\\0&1&0\end{pmatrix}.
\]
Its Smith factors are $(1,1,|n|)$ for arbitrary allowed
integers, independently of $r$.
Before attaching the cusp, the remaining degree-one group is
$\langle\gamma,g\mid3kg=n\gamma\rangle=\mathbb Z$.
Its cusp meridian is $-g$, a nonzero multiple $-n$ of a
primitive generator. Hence the relative map
$H_2(N,M)=\mathbb Z\to H_1(X^\circ)=\mathbb Z$ is injective,
so it contributes no extra free degree-two class.
We conclude $H_2(X)=\mathbb Z/|n|$.
The cusp has Euler number two and the other pieces have
Euler number zero. Duality and the universal coefficient
theorem give the remaining groups. When $|n|=1$, sphere
recognition applies exactly as in D.2.6.

The independent executable
\texttt{code/manifolds\_type\_c\_ii.py} checks the Weierstrass
identity, section, $j$-function, integral monodromy conjugacy,
symbolic determinant, and 168 full augmented Smith forms.
The preceding proof has no upper bound on $k,n,r$.
The current-source search is recorded in the report; this
statement makes no priority claim.
\end{proof}
